\section{Introduction}
\label{sec:introduction-v132}

A degeneracy locus usually forgets the map that defines it.  We study
an inverse problem in which neither the ambient embedding nor the
normal directions are retained: can the abstract nonreduced failure
scheme recover a quadratic pencil?  The answer is affirmative for
every complex pencil, including singular pencils.  The first relation, rather than its reduced determinant support,
carries this information.  The inverse is also global: a thickening over a projective reduction
recovers an entire varying pencil subbundle and its actual source
bundle, up to one constant left transformation.  No fibrewise choice
of coordinates or projective line-bundle ambiguity is left.  A
single unmarked local algebra suffices for the constant-pencil
problem; a Schubert line gives another geometric realization.

Let \(V\) be a complex vector space of dimension \(n\ge3\), and put
\[
 W=\Sym^2V,\qquad N=\binom{n+1}{2},\qquad
 R\in\Gr(2,W),\qquad S_R=W/R,\qquad p=N-2.
\]
The cube-zero algebra \(A_R=\C\oplus V\oplus S_R\) has multiplication
\(\Sym^2V\twoheadrightarrow S_R\), with
\((V\oplus S_R)S_R=0\).  On \(X_R=\Gr(n,V\oplus S_R)\), let
\(\mathcal K\) denote the tautological bundle.  For any \(m\ge2\), set
\begin{equation}\label{eq:failure-intro}
 \widehat D_R=V\!\left(\Fitt_0\coker[
 \Sym^m(\OO\oplus\mathcal K)\longrightarrow A_R\otimes\OO]\right).
\end{equation}
This is the maximal-minor scheme of multiplication, not merely its
support.  It is independent of \(m\ge2\).  Its reduction is the
Schubert divisor where \(\mathcal K\to V\) drops rank.  Write
\(B_R=\Gr(n,S_R)\), let \(\mathfrak b_R\) be its ideal in
\(\widehat D_R\), and define
\[
 Z_R^{[k]}=V(\mathfrak b_R^{k+1})\subset\widehat D_R,
       \qquad d=n+2p=n^2+2n-4.
\]
An isomorphism of these neighbourhoods means an abstract complex
scheme isomorphism; it is supplied with no ambient or tensor marking.

\begin{theorem}[Sharp intrinsic reconstruction]
\label{thm:principal-finite-v143}
For all quadratic pencils \(R,R'\subset\Sym^2V\),
\[
 Z_R^{[d]}\simeq Z_{R'}^{[d]}
       \quad\Longleftrightarrow\quad
 R'=gR\quad\text{for some }g\in\PGL(V).
\]
For every \(0\le k<d\), the isomorphism class of \(Z_R^{[k]}\)
is independent of \(R\).  Thus \(d\) is the smallest uniform
reconstruction order.  The same statements hold after restriction,
using the graph projection, to any Schubert line in \(B_R\).
The resulting abstract curve-supported neighbourhood is independent
of the choice of that line.
\end{theorem}

The Grassmannian statement is Theorem~\ref{thm:sharp-finite-pencil}.
The curve-supported statement is
Corollary~\ref{cor:schubert-line-reconstruction-v145}.  Its more general
form, Theorem~\ref{thm:curve-reconstruction-v145}, applies over any
smooth connected projective curve with a rank-\(n\) bundle having
nontrivial projectivization.  For example,
\(\OO_C(-P)\oplus\OO_C^{n-1}\) works on every such pointed curve.
The curve is intrinsic as the reduction; a chosen embedding of the
curve is not part of the isomorphism data.

\subsection{A moving family, not only its individual fibres}
The same finite object also reconstructs a pencil that varies over its
projective reduction.  Let \(B\) be a smooth connected projective
complex variety, let \(\mathcal U\) have rank \(n\), and let \(\mathcal R\) be a rank-two subbundle of
\(\Sym^2V\otimes\OO_B\) with locally free quotient.  Use the local
ideal \((\det T)I_p(\gamma\Sym^2T)\) in
\(\Hom(\mathcal U,V)\) and truncate at order \(d\).
Theorem~\ref{thm:unrestricted-moving-v146} states that two resulting
abstract finite schemes are isomorphic exactly when their data are
related by an isomorphism of reduced bases, one constant linear left
map, and an actual isomorphism of the source bundles, carrying one
pencil subbundle to the other.  The projection and both tensor factors
are forgotten in the input.  The varying coefficient line is recovered
as a subbundle, not by pointwise choices of projective coordinates.

This strengthens the constant-pencil curve theorem in two directions:
the reduced base and source bundle may vary in the comparison, and the
pencil need not be constant on that base.  Example~\ref{ex:nonconstant-pencil-v146}
has a nonconstant unordered cross-ratio and six spectral collision
points, all retained by the same finite scheme.  Theorem~\ref{thm:automorphism-kernel-v146}
then identifies the remaining freedom in an isomorphism.  Its quotient
is the group of the recovered linear data; its kernel acts trivially
on the associated graded scheme and has a terminating commutator
filtration.  Explicit nonlinear shears show why that kernel cannot be
omitted.  These statements are consequences of the unmarked inverse,
not additional claims of novelty for the classical Pluecker map or
for spectral cross-ratios.

\subsection{A single unmarked local algebra suffices}
The strongest form of the inverse needs no positive-dimensional support.
Choose an \(n\)-dimensional source space \(U\), write \(T=(t_{ij})\)
for its universal map to \(V\), and let \(\mathfrak m=(t_{ij})\).
Define
\[
 \mathfrak A_R^{[d]}=
 \C[t_{ij}]/\bigl((\det T)I_p(\gamma_R\Sym^2T)
                                      +\mathfrak m^{d+1}\bigr).
\]
\begin{theorem}[Local algebraic inverse]\label{thm:principal-local-v146}
For every pair of complex quadratic pencils, an ungraded algebra
isomorphism \(\mathfrak A_R^{[d]}\simeq\mathfrak A_{R'}^{[d]}\)
exists if and only if \(R'=gR\) for some \(g\in\PGL(V)\).
The order \(d=n^2+2n-4\) is the smallest uniform order already for
these single-point objects.
\end{theorem}
The full statement and proof are
Theorem~\ref{thm:artin-local-inverse-v146}.  Proposition~\ref{prop:first-relation-orbit-v147}
separates the general truncation argument from the pencil-specific
orbit rigidity: nonlinear coordinate changes form an explicitly described
unipotent kernel, whereas the first relation space carries the
reconstruction data.  The extra local ingredient
is that the residual coefficient module has factor-support ranks
\((1,\binom N2)\).  Transposition reverses these unequal ranks and
therefore cannot carry one pencil ideal to another.  This orients the
unordered matrix factors without a nontrivial projective bundle.
Theorem~\ref{thm:unrestricted-moving-v146} consequently removes that
restriction from the moving-family theorem as well.  This is not a
claim that the reduced determinant cone alone remembers the pencil:
the nonreduced first relation kernel and its residual coefficients
are indispensable.  Proposition~\ref{prop:local-automorphisms-v146}
also identifies the entire tangent-identity kernel, including its
affine-space parametrization and nonadditive group law.

\subsection{The intrinsic mechanism}
In the graph neighbourhood of \(B_R\), with tautological bundle
\(\mathcal U\) and graph map \(T:\mathcal U\to V\), the ideal is
\[
                 (\det T)I_p(\gamma_R\Sym^2T).
\]
Its generators have degree \(d\).  This degree calculation alone
only explains why lower orders are constant.  The inverse requires
three further intrinsic operations.  Multiplication in the finite
scheme recovers the first relation kernel and hence the whole
homogeneous cone.  The rank-one Fano scheme of its reduced determinant
cone then recovers the two tensor rulings.  On the socle Grassmannian
one may distinguish them by projective triviality, already on a
Schubert line.  In general the unequal coefficient-support ranks
orient them without that restriction.  Properness forces the induced
map to \(\PGL(V)\) to be constant.  After choosing its linear lift,
the actual normal-bundle map lies in the commutant of \(\End(V)\),
which recovers the right bundle without a line twist
(Lemma~\ref{lem:actual-factor-descent-v147}).  Finally, cancellation of the intrinsic
determinant ideal recovers the coefficient module.  For pencils its
exterior representation is irreducible, and exterior duality identifies
its single coefficient line with the Pluecker line of \(R\).

Theorem~\ref{thm:structural-coefficient-reconstruction} separates this
geometric mechanism from the particular multiplication algebra.  The
curve theorem shows why the mechanism is not tied to the dimension of
the original parameter Grassmannian.  The geometric orientation used here has a local replacement: coefficient-support asymmetry.  Theorem~\ref{thm:artin-local-inverse-v146} uses it to reconstruct the pencil from one unmarked Artin local algebra.

Theorem~\ref{thm:general-moving-coefficients-v147} gives the broader
principle.  It reconstructs simultaneously varying subspaces in
several Schur modules from the unmarked first-relation thickening,
provided one coefficient support is nonzero and proper.  The
coefficient spaces may have arbitrary prescribed ranks; they need
not be lines coming from quadratic pencils.  Its proof isolates the
actual factor-descent step from classical linear-preserver theory.
Example~\ref{ex:isotrivial-covers-v147} shows why the global descent is
not equivalent to fibrewise reconstruction: two Zariski locally
trivial families have the same Artin fibre and isomorphic nilradical
graded vector bundles in every degree, but their abstract total
schemes differ.  The inverse detects the inequivalent degree-three
coefficient maps.  This is the mechanism behind the moving-pencil
theorem, rather than
an assertion of novelty for canonical grading or for the numerical
value of the first relation degree.


\subsection{An exact image criterion and a pencil-native global consequence}
The coefficient inverse admits an intrinsic converse.
Theorem~\ref{thm:recognition-dichotomy-v148} starts with an
unmarked first relation whose homogeneous radical is a determinant
cone.  After recovering its unordered rulings, invariance under one
full factor group is necessary and sufficient for the residual
relation to come from coefficient subspaces in the Cauchy modules.
Both orderings satisfy this test precisely when every coefficient
support is zero or full.  Thus the strict-support condition is
exactly the criterion for a unique admissible orientation in the
recognized class.  Theorem~\ref{thm:exact-coefficient-stabilizer-v148}
computes its entire affine stabilizer, including the transpose
component and the universal nonlinear automorphism kernel.
Corollary~\ref{cor:intrinsic-orbits-v148} expresses the inverse as
an equivalence of geometric orbit sets, with those ambiguities
specified rather than suppressed.

A global consequence lies inside the quadratic-pencil problem
itself.  Fix a plane \(H=\langle x,y\rangle\subset V\).
For a degree-\(m\) map \(f=[a:b]:\Pj^1\to\Pj(H)\), take
the pencil subbundle
\[
                  \mathcal R_f=(ax+by)H.
\]
Every fibre is equivalent to \(\langle x^2,xy\rangle\).
Nevertheless, Theorem~\ref{thm:covering-pencil-moduli-v148} proves
\[
 Y_f\simeq Y_{f'}\quad\Longleftrightarrow\quad
 f'=\alpha\circ f\circ\sigma,\qquad
                         \alpha,\sigma\in\PGL_2(\C),
\]
for their unmarked order-\(d\) failure thickenings with trivial
source bundle.  Their fibre algebras and every nilradical graded
vector bundle depend only on \(n,m\), not on \(f\).
Their abstract pencil, quotient and first-relation bundles are
fixed as well.  Thus the entire covering equivalence problem, not
only one chosen pair, is retained in the global multiplication and
is invisible to those additive data.  The maps \(z^3\) and
\(z^3+z\) give an explicit separating pair of quadratic pencils
(Corollary~\ref{cor:pencil-native-pair-v148}).

Neither complete reducibility of Cauchy modules, linear preservers
of a Segre cone, nor division by the common factor of a pencil is
new in isolation.  Their synthesis here has a precise output:
recognition of the image and its exact orientation ambiguity after
all matrix markings are forgotten, followed by recovery of a
whole covering map from a locally isotrivial pencil thickening.
A projective factor map alone would not recover the actual source
bundle; the intrinsic normal map and its commutant do so.  A
pointwise pencil inverse would not recover a covering map; one
constant left transformation over the entire proper base does so.
These distinctions specify the geometric content beyond the
classical ingredients.  The recognized coefficient class is not the
class of all determinant-radical relations.  Its geometric orbit
statement alone does not describe nonreduced deformation groupoids;
the relative comparison below uses a further scheme-theoretic inverse.


\subsection{Intrinsic strata, infinitesimal parameters, and pencil moduli}
The relative inverse begins with a geometric condition on the first
relation, not invariance under a factor group.  For a space \(K\) of
homogeneous forms, prescribe the degree of its greatest common divisor
and give that locus its reduced locally closed scheme structure.
Lemma~\ref{lem:universal-gcd} proves that multiplication by the
universal common-divisor line is an isomorphism onto this stratum.
Consequently the common factor is recovered even over a nonreduced
parameter base.  The condition is scheme-theoretic: merely asking all
geometric fibres to have the specified common factor is insufficient.
When that divisor is a power of a determinant, the remaining
parameter is an arbitrary gcd-free residual subspace, modulo the
full determinant-preserver group
(Theorem~\ref{thm:determinant-stratum}).  This is a smooth quotient
stack, and neither of the two factor symmetries is required.

For pencils the exact common divisor is \(\delta^{n-1}\), and the
primitive residual relation has degree \(3n-4\)
(Lemma~\ref{lem:primitive-pencil}).  The Cauchy projection and the
Pluecker equations recognize its image as a locally closed scheme,
not just a set of complex points.  The resulting relative theorem is
\begin{equation}\label{eq:intro-effective-moduli}
 \mathscr F_{\rm pen}^{\rm eff}\simeq
                [\Gr(2,\Sym^2V)/\PGL(V)].
\end{equation}
Here the superscript records two explicit operations on morphisms:
first remove tangent-identity nonlinear substitutions, then remove
the ineffective right matrix factor, and fppf stackify.
Theorem~\ref{thm:pencil-stack-equivalence} proves this equivalence
over arbitrary complex parameter schemes, retaining all singular
pencils.  It does not identify the unrigidified algebra stack with
the pencil stack.  Its tangent complex is
\([\mathfrak{pgl}(V)\to\Hom(R,\Sym^2V/R)]\), so the effective
stack is smooth of dimension \(n-3\), with precisely the pencil
stabilizer jumps.  Rank-incidence and spectral Fitting strata retain
their scheme structures and deformation equations by
Corollary~\ref{cor:pencil-specialization-stack}.  This comparison
concerns the natural parameter space of all pencils, rather than an
external moduli problem inserted into a common-factor example.

The ambient determinant-divisor stratum is substantially larger.
Theorem~\ref{thm:transverse-pencil-directions} computes its normal
space to the pencil stratum in three parts: non-Pluecker lines,
non-scalar right-factor directions, and other Cauchy summands.
Theorem~\ref{thm:fixed-support-grassmannians} constructs through every
pencil a projective Grassmannian of relations with the same determinant
radical, containing a projective line whose other points admit neither
factor symmetry.  Thus the family inverse describes how the
coefficient and pencil loci sit in, and can be left within, a larger
intrinsic stratum.  It does not falsely infer deformation stability
from fibrewise recognition.

\subsection{Framed families and spectral consequences}
Once tensor coordinates are retained, the parameter map is simpler.
The first relation spaces give a closed immersion of
\(\Gr(2,W)\) into a Grassmannian, by
Theorem~\ref{thm:finite-parameter-immersion}.  Its proof factors through
the classical Pluecker embedding, tensoring a line with a fixed space,
and a fixed linear inclusion.  This is the framed encoding consequence
of the coefficient calculation, not a substitute for the unmarked
inverse.  Theorem~\ref{thm:universal-finite-neighbourhood} glues the
actual universal order-\(d\) neighbourhood.  It is finite locally
free over the relative socle of rank
\(\binom{n^2+d}{d}-\binom N2\), and projective and flat over the
pencil parameter scheme.  It commutes with arbitrary complex base
change with that socle retained.  Over a nonreduced base, the specified
relative socle ideal need not be the absolute nilradical.  These marked statements alone do not identify unmarked moduli
stacks or deformation groupoids.  The common-divisor descent and
ineffective-inertia analysis in
Theorem~\ref{thm:pencil-stack-equivalence} supply that separate step.

For regular pencils, the spectral sheaf and its Fitting incidence
schemes are read from the reconstructed pencil.  We retain these
consequences because they exhibit information genuinely lost by
coarser invariants: fixed discriminants and fixed reduced rank loci
can conceal different elementary divisors.  In
Theorem~\ref{thm:fixed-spectral-classification}, every classical
dominance specialization is transported to a flat family of the
finite neighbourhoods, with every strict transition detected at
order \(d\).  The spectral orbit classification itself is classical,
not a second independent Torelli theorem.

\subsection{Prior work and the publication object}
The Weierstrass--Segre classification is used in its root-labelled
form; compare \cite[Theorem~1.1 and Corollary~2.1]{FevolaMandelshtamSturmfels}.
The orthogonal nilpotent input is isolated in
Lemma~\ref{lem:orthogonal-orbits-v145}, with precise references to
\cite[Proposition~1 and Theorem~1]{Ohta1986}, an independent dimension
calculation, and an explicit distinction between full orthogonal
orbits and their special-orthogonal components.

Ballico's relative multiplication construction
\cite[Theorem~0.2; \S1, Remark~1.2 and equation~(6)]{Ballico96}
is a substantive antecedent for studying quadratic-normality failure
on finite linear sections of an embedded surface.  Here the defining
algebra varies, and the inverse starts from an unmarked infinitesimal
failure scheme.  The opening page of the distinct earlier paper \cite[p.~5]{Ballico93}
concerns higher-order embedding properties and restriction maps to
finite subschemes.  Its remaining theorem and proof text has not
been inspected; that first page does not establish the absence of
an inverse statement.  No anticipation or nonanticipation conclusion
relative to that article is asserted.

The Artin-local aspect is compared at theorem level with
\cite{EliasRossiShort,EliasRossiCompressed} in
Section~\ref{sec:first-relations-v147}, including the distinction
between canonical grading of a nonhomogeneous local algebra and
orbit rigidity of our already homogeneous first relation.  The
linear-preserver input is classical \cite{MarcusMoyls}; its precise
hypotheses and a proof in the required case are included in
Lemma~\ref{lem:linear-preserver-v147}.  The algebraic-group
quotient step is supplied by the homomorphism theorem
\cite[Theorem~5.39 and Remark~5.42]{Milne2017}; its use in the
exact stabilizer calculation distinguishes a scheme-theoretic
quotient from a bare bijection of complex points.

The present article consists of intrinsic reconstruction, its
relative deformation comparison, and its spectral readout.  The reciprocal-likelihood theorem and projective critical
schemes are preserved in a separate applications manuscript.  Its
finite-flat critical divisor requires supplied split semisimple data
and simple disjoint projective poles; its totally real cover requires
a supplied real definite realization and separated data.  None of
those hypotheses enters the all-pencil theorem here.  Independent
web, exterior-operator, K3, polarized, and boundary results are
preserved in an explicitly non-submitted research archive.  They are
not inputs to any proof in this article and are not part of its
significance claim.  Every proof required here is included here.
